\section{Codebase layout and terminology}
\label{app:codebase}

The analysis is a chain of four components; the italicized names are used
throughout the note.

\begin{itemize}[itemsep=2pt]
  \item \emph{Event loop} (C++): \texttt{runEventLoopOmniFold}, built on the
        MINERvA Analysis Toolkit (MAT) and the MINERvA-101 framework. It
        reads the production AnaTuples, applies the selection and the
        systematic-universe machinery, and writes per-event ROOT trees.
  \item \emph{Omnifiles}: the event-loop outputs --- per-event records for
        data, truth/reco-paired simulated signal, simulated background, and
        the truth-level efficiency denominator, with POT bookkeeping; the
        \texttt{\_universes\_full} variants add per-universe weight and
        shifted-kinematics branches.
  \item \emph{Driver} (Python): the per-analysis orchestration ---
        \texttt{2d-unfolding/unfold\_2d\_omnifold\_unbinned.py} and its
        $N$-dimensional generalization
        \texttt{nd-unfolding/unfold\_nd\_omnifold\_unbinned.py}. The driver
        assembles training arrays from the omnifiles, applies the
        purity-weight background subtraction (\S\ref{sec:method-algo}),
        invokes the engine, and extracts the cross section.
  \item \emph{Engine} (Python):
        \texttt{unbinned\_unfolding/python/omnifold.py} --- the OmniFold
        iteration itself, agnostic to the learner backend (LightGBM in
        production; sklearn \texttt{exact}/\texttt{hist} alternates; MLP and
        point-cloud-transformer cross-checks).
\end{itemize}

The recoil-point-cloud cross-checks of \S\ref{sec:pet} replace the driver's
scalar classifier features with the non-muon cluster and truth-hadron inputs
of \texttt{nd-unfolding/pet/}, reusing the same omnifile and cross-section
machinery. Muon scalars are retained for selection and binning but are not
classifier inputs in these legacy PET products.
